\documentclass[11pt,a4paper]{article}
\usepackage[T1]{fontenc}
\usepackage[utf8]{inputenc}
\usepackage{lmodern}
\usepackage[margin=25mm]{geometry}
\usepackage{amsmath,amssymb,amsthm,mathtools,booktabs,array}
\usepackage{microtype}
\usepackage[hyphens]{url}
\usepackage[colorlinks=true,linkcolor=blue,urlcolor=blue,citecolor=blue,hypertexnames=false]{hyperref}
\usepackage{enumitem}
\hypersetup{pdftitle={Erdos Problem 708: Verified Partial Arguments and the Remaining Gap},pdfauthor={Working research note prepared for Denis Paliy},pdfsubject={Partial research results; the general conjecture remains unresolved}}
\setlist{itemsep=3pt,topsep=5pt}
\newtheorem{theorem}{Theorem}[section]
\newtheorem{lemma}[theorem]{Lemma}
\newtheorem{proposition}[theorem]{Proposition}
\newtheorem{corollary}[theorem]{Corollary}
\theoremstyle{remark}
\newtheorem*{remark}{Scope}
\DeclareMathOperator{\lcm}{lcm}
\newcommand{\vp}{v_p}
\newcommand{\pos}[1]{\left(#1\right)_+}
\newcommand{\PP}{\mathcal P}
\newcommand{\II}{I_{x,M}}
\title{Erd\H{o}s Problem 708:\\Verified Partial Arguments and the Remaining Gap}
\author{Working research note prepared for Denis Paliy}
\date{22 November 2026}
\begin{document}
\maketitle
\begin{center}
\fbox{\parbox{0.91\textwidth}{\textbf{Status: the general problem is not solved in this note.}
Neither $g(n)\le 2n$ nor $g(n)\le(2+o(1))n$ has been proved here.
Theorems explicitly marked conditional depend on the unproved inequality
$(H_2)$. The finite computations below do not establish a universal theorem.
This is an AI-assisted research record, with no claim of literature priority,
formal verification, journal acceptance, or independent human review.}}
\end{center}

\begin{abstract}
We give a self-contained whole-cofactor extension argument valid at every
interval length, including a saving for each distinct largest prime among
the nonsmooth sources. Combined with a boxed linear relaxation, it gives a
conditional bound $g(n)\le 2n+\pi(q_n)-1$, where
$q_n=\min\{Z\ge2:Z\ge2n+\pi(Z)-1\}$, from the global hinge inequality
$(H_2)$. Thus this conditional route has no remaining interval-size gap.
An unconditional consequence covers sets whose every member has a prime
factor greater than $2n-1$. We also give exact obstructions to two
overly strong proof strategies, including a family on which a particular
independent-prime charging bound has unbounded relative error while a
valid cover uses $2n-2$ positions. Proofs, computations, assumptions, and
attribution are separated throughout.
\end{abstract}

\section{The intended problem and the proof boundary}
Let $A\subset\mathbb Z_{\ge2}$ have $n$ distinct elements. Put
$M=\max A$ and
\[
 \II=\{x+1,\ldots,x+M\},\qquad x\in\mathbb Z_{\ge0}.
\]
Define
\[
 \tau(A,\II)=\min\left\{|B|:B\subseteq\II,\quad
              \prod_{a\in A}a\mid\prod_{b\in B}b\right\},
 \qquad g(n)=\sup_{A,x}\tau(A,\II).
\]
Every chosen interval position is used at most once. We use the intended
``at most'' convention, not a literal exact-cardinality wording that can
fail merely because an interval is too short. The questions are
\[
 g(n)\stackrel{?}{\le}2n,\qquad
 g(n)\stackrel{?}{\le}(2+o(1))n.
\]
The recent finite-case paper of Tripodo explicitly leaves the general
question open \cite{tripodo}. This note addresses that general question,
without presenting small cases or a conditional reduction as its solution.

Several arguments below work more generally with $M\ge\max A$.
The whole interval is always a cover, because
\[
 \prod_{a\in A}a\mid M!,\qquad
 \prod_{b\in\II}b=M!\binom{x+M}{M}.
\]
Consequently $M\le2n$ is immediate. Intervals containing zero are trivial;
all-negative intervals reduce to positive intervals by reflection.

\paragraph{Attribution.}
The linear-programming formulation and the global raw-relaxation transfer
are developed in Wang's source \cite{wang}. The cumulative-gcd extension
and proper-divisor matching approach occur in Bergstr\"om's manuscripts
\cite{berg-asym,berg-hinge}. The collision charging argument is already in
Chen and is restated in Preil's notes \cite{chen,preil}. We rederive the
parts used below and extend the whole-cofactor statement to all scales.
No assertion of novelty or priority is made for these reformulations.

\section{Whole-cofactor extension without a scale restriction}
The next elementary lemma will keep repeated prime powers properly paid.

\begin{lemma}[One-step extension]\label{lem:extension}
Suppose $H\mid\prod_{b\in B}b$, where $B\subseteq\II$,
$1\le h\le M$, and
\[
 \gcd(H,h)\le\left\lfloor\frac{M}{|B|+1}\right\rfloor.
\]
Then one unused position can be added so that the enlarged product is
divisible by $Hh$. A unit request $h=1$ can instead be skipped.
\end{lemma}
\begin{proof}
If an unused multiple of $h$ exists, select it. Otherwise at least one
multiple of $h$ lies in $B$, so $h$ and $H$ both divide the current product.
Therefore $\lcm(H,h)$ divides that product. Write $d=\gcd(H,h)$.
The interval has at least $\lfloor M/d\rfloor\ge |B|+1$ multiples of $d$;
one is unused. Multiplying by it supplies the missing factor in
$Hh=d\lcm(H,h)$.
\end{proof}

Fix an integer $Z\ge2$. Let $A_0$ be the sources with largest prime factor
at most $Z$. Write $r=|A\setminus A_0|$ and let $t$ be the number of
distinct largest prime factors of sources in $A\setminus A_0$.

\begin{theorem}[All-scale whole-cofactor extension]\label{thm:peeling}
Let $B_0\subseteq\II$ cover $A_0$, and suppose $r>0$. Set
\[
 K=|B_0|+2r-t.
\]
If $Z\ge K$, then $B_0$ extends to a cover of $A$ of size at most $K$.
There is no assumption $Z^2>M$.
\end{theorem}
\begin{proof}
Group the remaining sources by their largest prime factor $p>Z$:
$a_{p,1},\ldots,a_{p,m_p}$. We maintain a target integer $H$ known to
divide the product of the selected positions. Extra prime factors in
those positions do not change the definition of $H$.

\emph{Stage 1: one complete source per group.}
Process the groups in increasing $p$. Before a group is processed, $H$
is the product of the smooth sources and one source from each earlier
group. Thus $p\nmid H$, and
\[
 \gcd(H,a_{p,1})\le a_{p,1}/p\le M/p<M/Z.
\]
Before every addition $|B|+1\le |B_0|+t\le K\le Z$. The integer gcd
therefore satisfies Lemma~\ref{lem:extension}. This stage uses at most
$t$ additional positions and multiplies $H$ by each $a_{p,1}$.

\emph{Stage 2: the chosen prime of each remaining source.}
For each group we must multiply the target by $p^{m_p-1}$.
If $p^2\le M$, the interval contains at least
\[
 \lfloor M/p\rfloor\ge p>Z\ge K
\]
multiples of $p$. A fresh multiple exists before every requested
addition. Adding one per copy is valid even if other parts of the target
also contain $p$.

If $p^2>M$, every source contains $p$ to exponent at most one. Moreover
no source in a different group contains $p$: a larger largest prime $q$
would give $pq>M$, and a smaller largest prime cannot occur together
with $p$. Thus the target before processing this group contains exactly
one copy of $p$. The actual current product may contain more; write its
$p$-valuation as $u\ge1$. If $u\ge m_p$, nothing is needed. Otherwise
there are at least $\lfloor M/p\rfloor\ge m_p$ interval multiples of $p$,
because the $m_p$ sources are distinct multiples of $p$ at most $M$.
At most $u$ selected positions are divisible by $p$. Hence at least
$m_p-u$ are unused; selecting these suffices. This adds at most $m_p-1$
positions and preserves every other required valuation.
Summed over groups, Stage 2 uses at most $r-t$ new positions.

\emph{Stage 3: whole remaining cofactors.}
For each $j\ge2$ put $d=a_{p,j}/p$. If $d>1$, then
\[
 |\II\cap d\mathbb Z|\ge\lfloor M/d\rfloor\ge p>Z\ge K.
\]
Select a fresh multiple of each such $d$, disjoint from all earlier
positions. The cofactor may contain more copies of $p$ or other primes
larger than $Z$; all of it is paid on this fresh position. At most $r-t$
positions are added. Finally,
\[
 |B|\le |B_0|+t+(r-t)+(r-t)=K,
\]
and the factorization of the sources proves the required divisibility.
\end{proof}

\begin{corollary}[An unconditional class]\label{cor:rough}
For $n\ge2$, suppose each $a\in A$ has a prime divisor greater than
$2n-1$. If $t$ is the number of distinct largest prime factors in $A$,
then
\[
 \tau(A,\II)\le2n-t\le2n-1.
\]
\end{corollary}
\begin{proof}
Set $Z=2n-1$, $A_0=\varnothing$, $B_0=\varnothing$, and $r=n$.
Since $t\ge1$, the budget $K=2n-t$ is at most $Z$. Apply
Theorem~\ref{thm:peeling}.
\end{proof}
\begin{remark}
This corollary is a restricted theorem, not an assumption that can be
made for an arbitrary input. Integers composed of primes at most $2n-1$
remain a genuine obstruction to applying it to all of Problem 708.
\end{remark}

\section{The exact missing inequality and a conditional consequence}
For weights $w_p\in[0,1]$ on primes $p\le M$, set
\[
 X_w(k)=\sum_{p\le M}w_pv_p(k),\qquad u_+=\max(u,0).
\]
The unproved global inequality in this route is
\begin{equation}\tag{$H_2$}\label{eq:h2}
 \sum_{a=1}^{M}\pos{X_w(a)-2}
 \le \sum_{b\in\II}\pos{X_w(b)-1}
 \quad\text{for every }M,x,w.
\end{equation}
We do not prove \eqref{eq:h2} here. The following consequences keep it
as an explicit hypothesis.

\subsection{Boxed relaxation and dual clipping}
For $D_p(A)=\sum_{a\in A}v_p(a)$, define
\[
 L(A,I)=\min\left\{\sum_{b\in I}z_b:
 0\le z_b\le1,\quad\sum_{b\in I}v_p(b)z_b\ge D_p(A)
 \text{ for all demanded }p\right\}.
\]
The upper boxes $z_b\le1$ are essential: an interval position cannot be
used repeatedly. The term ``raw'' in the cited source refers to the
coefficients $v_p(b)$, not to removal of those boxes.
The feasible set is nonempty by the whole-interval observation. Finite
linear-programming duality gives
\begin{equation}\label{eq:dual}
 L(A,I)=\max_{w_p\ge0}\left[
  \sum_{a\in A}X_w(a)-\sum_{b\in I}\pos{X_w(b)-1}\right].
\end{equation}
If $w_p=1+\delta>1$, lower it to one. For every multiple of $p$ the
right hinge remains in its linear region. Its total decrease is exactly
$\delta\sum_{b\in I}v_p(b)$, whereas the source term decreases by
$\delta D_p(A)$. The former is at least the latter, because
\[
 \sum_{b\in I}v_p(b)
 \ge\sum_{e\ge1}\lfloor M/p^e\rfloor
 \ge D_p(A).
\]
Thus clipping cannot reduce the dual objective, and weights in $[0,1]$
suffice. If \eqref{eq:h2} holds, then
\[
 \sum_{a\in A}X_w(a)
 \le 2|A|+\sum_{a=1}^M\pos{X_w(a)-2}
 \le 2|A|+\sum_{b\in I}\pos{X_w(b)-1}.
\]
Therefore $L(A,I)\le2|A|$.

Conversely, if $L(A,I)\le2|A|$ holds for all source subsets of
$\{1,\ldots,M\}$, apply \eqref{eq:dual} to
$A=\{a\le M:X_w(a)>2\}$ to recover \eqref{eq:h2}.
The hinge formulation is therefore a substantive covering statement,
not an innocuous estimate that can be omitted from the proof.

\subsection{Rounding and all interval lengths}
\begin{theorem}[Conditional; assumes global $(H_2)$]\label{thm:conditional}
Let
\[
 q_n=\min\{Z\in\mathbb Z:Z\ge2,\quad Z\ge2n+\pi(Z)-1\}.
\]
If the global inequality \eqref{eq:h2} holds, then
\[
 g(n)\le 2n+\pi(q_n)-1
       =2n+(2+o(1))\frac{n}{\log n}.
\]
In particular the asymptotic question would be answered for every
interval length. This does not prove the exact inequality $g(n)\le2n$.
\end{theorem}
\begin{proof}
The minimum exists: $Z=4n$ works since $\pi(4n)\le2n$.
Split the sources at $Z=q_n$, and let $m$ be the number of smooth
sources. If $m>0$, solve their boxed relaxation. An optimal vertex has
at most $\pi(Z)$ fractional coordinates: otherwise its fractional
columns are dependent on the tight demand rows, permitting a nonzero
two-sided feasible perturbation. Rounding all fractional coordinates up
preserves feasibility. If some coordinate is fractional, the increase
is strictly less than their number; if none is fractional, the solution
is already integral. Since $L\le2m$ and $\pi(Z)\ge1$, both cases give
\[
 |B_0|\le2m+\pi(Z)-1.
\]
If $m=0$, take $B_0=\varnothing$ instead. If there are no rough sources,
we are done by the same rounding bound. Otherwise the budget in
Theorem~\ref{thm:peeling} satisfies
\[
 |B_0|+2(n-m)-t\le2n+\pi(Z)-1\le Z.
\]
For $m=0$ it is even at most $2n-1$. The extension theorem applies at
every $M$ and gives the finite bound. The integer function $Z-\pi(Z)$
advances by zero or one at each step, so minimality (including the
immediate case $n=1$) gives $q_n-\pi(q_n)=2n-1$.
Together with $q_n\le4n$ and $\pi(x)=O(x/\log x)$, this implies
$q_n=2n+O(n/\log n)$; the prime number
theorem yields $\pi(q_n)=(2+o(1))n/\log n$.
\end{proof}

Wang's existing raw transfer already implies the sufficient conditional
bound $g(n)\le2n+\pi(4n)-1$ from $(H_2)$ for all $M$ \cite{wang}.
The threshold refinement above does not supply the missing hypothesis.

\section{A collision charging bound cannot close the problem by itself}
Choose, for each $a\in A$, an interval multiple $b_a$. Let $t$ be the
number of distinct chosen multiples and let $S_b$ be the union of the
prime supports of sources assigned to $b$. The known collision bound is
\begin{equation}\label{eq:hub}
 \tau(A,I)\le U:=t+\sum_{a\in A}\omega(a)-\sum_b|S_b|,
\end{equation}
where $\omega$ counts distinct prime factors \cite{chen,preil}.
For completeness, for each prime $p$ the source-to-target graph matching
$a$ to multiples of $p^{v_p(a)}$ satisfies Hall by nested neighborhoods:
\[
 \#\{a:p^e\mid a\}\le\lfloor M/p^e\rfloor
 \le\#\{b\in I:p^e\mid b\}.
\]
Start a partial matching by assigning the largest $p$-valuation source
in each relevant collision group to its chosen multiple. Augmenting
paths extend it to a full matching while preserving the set of matched
right vertices, although individual assignments can change. If $k_p$
sources contain $p$ and $c_p$ such groups exist, this adds $k_p-c_p$
new right vertices. Taking their union over $p$ proves \eqref{eq:hub}.

\subsection{A small exact obstruction}
Take
\[
 A=\{60,66,70,78,84,90\},\quad M=90,\quad L=\lcm(A)=180180,
\]
and $I=\{L-44,\ldots,L+45\}$. Each source exceeds $45$ and divides
$L$, so its unique interval multiple is $L$. Every source has three
distinct prime factors, whose union is $\{2,3,5,7,11,13\}$. Therefore
every choice forces
\[
 U=1+18-6=13>12=2n.
\]
The actual product-cover optimum is only four, as established by exact
0/1 dynamic programming in the accompanying code. A cover is
\[
 B=\{180140,180144,180145,180180\}.
\]
Its valuations at the demanded primes are
\[
\begin{array}{c|rrrrrr}
 &2&3&5&7&11&13\\\hline
 \prod A&8&6&3&2&1&1\\
 \prod B&8&6&3&2&1&1
\end{array}
\]
The displayed cover alone suffices to show that this obstruction is to
the charging bound, not to Problem 708. The minimality claim is separately
certified by exhaustive dynamic programming, not by a solver tolerance.

\subsection{Unbounded overestimate on a whole family}
\begin{proposition}\label{prop:family}
For any fixed integer $s\ge3$, there is a family of genuine instances
with $\gcd(A)=1$, connected prime support, and $2n<M<16n^2$, on which
the forced expression $U$ in \eqref{eq:hub} is at least $(s-o(1))n$.
For the same intervals there is an explicit cover with $2n-2$ positions.
\end{proposition}
\begin{proof}
Let $D$ be the product of the first $s$ primes. Let odd $K\to\infty$
satisfy $\gcd(K,D)=1$ and $K\ge2D$. Set
\[
 A_K=\{Dk:(K+1)/2\le k\le K\}\cup\{(D-1)K\},
 \quad n=(K+3)/2,\quad M=DK.
\]
The extra source is not divisible by $D$ and is distinct from the others.
The consecutive indices in the first group have gcd one, so
$\gcd(A_K)=\gcd(D,(D-1)K)=1$. The first group shares the prime support
of $D$, and the extra source shares the prime factors of $K$ with $DK$;
hence the support is connected. All sources exceed $M/2$.

Put $L=\lcm(A_K)$ and take
$I_K=\{L-M/2+1,\ldots,L+M/2\}$. Every source has the unique multiple
$L$. All active primes are at most $K$, and each source in the first
group has at least $s$ distinct prime factors. Thus
\[
 U=1+\sum_{a\in A_K}\omega(a)-|\PP(A_K)|
 \ge1+s(n-1)-\pi(K)=(s-o(1))n.
\]
The size conditions on $M$ hold for all sufficiently large $K$.

To construct a small actual cover, select all $K$ multiples of $D$ in
$I_K$. After division by $D$, they are $K$ consecutive integers whose
product is divisible by $K!$. The product of the original selected
multiples is therefore divisible by $D^K K!$. Writing $h=(K-1)/2$, its
product divided by the first group's source product is a multiple of
$D^h h!$. As $h\ge D-1$, this supplies $D-1$. There are exactly $D$
multiples of $K$ in $I_K$, and exactly one is also a multiple of $D$
because $\gcd(K,D)=1$. Add one of the others; it supplies the remaining
factor $K$ without reusing a position. The cover size is
$K+1=2n-2$.
\end{proof}

This family rules out the proposed universal numerical bound $U\le2n$
even after simple gcd and connectedness reductions. It does not rule out
using the same initial multiples with a more effective joint repair.

\section{Proper-divisor matching: a sufficient hypothesis, not a proof}
Let $\Omega(m)=\sum_p v_p(m)$ count prime factors with multiplicity, and
let $\mathcal S_M=\{m\le M:\Omega(m)\ge3\}$. The following sufficient
condition appears in Bergstr\"om's work \cite{berg-asym,berg-hinge}:
\begin{equation}\label{eq:matching}
 \text{There is an injection }f:\mathcal S_M\hookrightarrow\II
 \text{ with }(m/p)\mid f(m)\text{ for some prime }p\mid m.
\end{equation}
If it holds, then $X_w(f(m))\ge X_w(m)-1$. Since $X_w(m)>2$ requires
$\Omega(m)\ge3$, summing
$\pos{X_w(f(m))-1}\ge\pos{X_w(m)-2}$ proves $(H_2)$.
No universal proof of \eqref{eq:matching} was obtained in this work.

\subsection{Removing only the largest prime is false}
Here is an exact small obstruction to that restricted rule. Set
\[
 M=440,\qquad x=4400,\qquad I=\{4401,\ldots,4840\},
\]
and let
\[
 S=\{2k:4\le k\le220,\ k\text{ composite},\ k\ne169\}.
\]
There are $220-1-\pi(220)-1=171$ sources in $S$, all with
$\Omega(m)\ge3$. The neighbor set for the rule
$(m/P^+(m))\mid b$ is exactly the set of interval multiples of
\[
 4,\ 6,\ 10,\ 14,\ 22.
\]
Indeed every composite $k\le220$ other than $13^2$ has a prime factor
at most $11$ remaining after its largest prime is removed. Conversely
the sources $2p^2$, $p\in\{2,3,5,7,11\}$, generate all five displayed
divisibility classes.

To count the neighbors, write even positions as $4620+2j$ with
$-109\le j\le110$. Because $2310=2\cdot3\cdot5\cdot7\cdot11$,
a position is a neighbor exactly when $j$ is divisible by one of these
five primes. There are $85$ such positive $j\le110$, $84$ negative
ones, and zero, for a total of $170$. The positive excluded integers
are $1$ and the $24$ primes from $13$ to $109$.
Thus Hall fails: $171>170$. Allowing removal of any prime, as in
\eqref{eq:matching}, is a different assertion and is not refuted.

\subsection{The full rule also fails if all composites are included}
Replacing $\Omega(m)\ge3$ by $\Omega(m)\ge2$ makes the full any-prime
rule false. An exact certificate in the archive uses
\[
 M=11763,\qquad H=5881,\qquad \pi(M)=1409.
\]
It supplies a set $T$ of $1411$ distinct offsets in $[-H,H]$, together
with a residue $c_p$ for every prime $p\le M$, satisfying
\[
 c_p+k\not\equiv0\pmod p\quad(k\in T).
\]
The Chinese remainder theorem gives a center $C>H$ with
$C\equiv c_p\pmod p$ for every such prime. In the $M$-term interval
$[C-H,C+H]$, all $1411$ positions $C+k$ have no prime factor at most
$M$. None can be a neighbor of a composite source $m\le M$, since
every quotient $m/p$ is a nontrivial integer at most $M$.

There are $M-1-\pi(M)=10353$ composite sources but at most
$M-1411=10352$ neighbors, proving a Hall violation. The exact verifier
checks every residue condition, a hexadecimal CRT representative, and
the whole neighborhood. For all-composite sources that neighborhood is
exactly the interval positions divisible by a prime $q\le M/2$:
every nontrivial quotient has such a prime divisor, while source $2q$
permits quotient $q$. Its verified size is $10352$.

This construction assumes no prime-tuple conjecture. The certificate
and its standard-library verifier are the finite evidence for the
displayed counts; no minimality or priority is claimed. It refutes only
the extension to $\Omega\ge2$, not \eqref{eq:matching} with
$\Omega\ge3$, $(H_2)$, or Problem 708. In fact this all-composite
instance has $M<2|A|$, so the original $2n$ cover bound is immediate by
selecting the whole interval.

\subsection{Finite experiments and their limits}
The matching code independently checks every returned assignment by
integer divisibility. A failed matching exports a Hall-deficient set
and recomputes its complete neighborhood exactly.
The unrestricted rule was tested in the following two batches:
\begin{center}
\begin{tabular}{@{}lll@{}}\toprule
Lengths $M$ & Shifts per length & Executed checks\\\midrule
Every $8\le M\le250$ & 100 & 24,300\\
$M=10,20,\ldots,2000$ & 100 & 20,000\\\bottomrule
\end{tabular}
\end{center}
No failure was found. These are sampled shifts, not exhaustive residue
classes; some checks overlap between batches. Shifts include small
integers, $M^2$, $M!-\lfloor M/2\rfloor$, and seeded random large
integers. The results do not prove \eqref{eq:matching}, $(H_2)$, or either
version of the original conjecture.

\section{Two analytical shortcuts that do not discharge \texorpdfstring{$(H_2)$}{(H2)}}
\subsection{A false sparse numerical comparison and a valid replacement}
The sparse binary argument in \cite{berg-hinge} uses the comparison
\[
 \frac{\mu}{1+\mu}>\frac{\mu^2+\sigma_2}{2},
 \qquad \mu=\sum_{p\in S}\frac1p<1,
 \quad \sigma_2=\sum_{p\in S}\frac1{p^2}.
\]
It is false for $S=\{2,3,7\}$:
\[
 \mu=\frac{41}{42},\quad \sigma_2=\frac{673}{1764},\qquad
 \frac{41}{83}<\frac{1177}{1764}.
\]
This refutes that numerical step, not $(H_2)$.

A correct elementary substitute is available with an explicit condition.
For $k=|S|$ and any interval $J$ of $L$ positive integers, two-term
Bonferroni gives
\[
 \#\{b\in J:\exists p\in S,\ p\mid b\}
 \ge L\left(\mu-\frac{\mu^2-\sigma_2}{2}\right)
      -\frac{k(k+1)}2.
\]
Meanwhile a union bound over prime squares and distinct prime pairs gives
\[
 \#\{a\le L:\Omega_S(a)\ge2\}
 \le\frac L2(\mu^2+\sigma_2).
\]
Thus the first count dominates the second whenever
$L\mu(1-\mu)\ge k(k+1)/2$. This is only a level-count comparison under
an additional hypothesis; it is not the full fractional-weight theorem.

\subsection{Averaging binary weights loses a linear term}
For $t\in[0,1]$, let $S_t=\{p:w_p\ge t\}$ and
$\Omega_t(b)=\sum_{p\in S_t}v_p(b)$. Then
\[
 X_w(b)=\int_0^1\Omega_t(b)\,dt,\qquad
 \int_0^1\pos{\Omega_t(b)-1}\,dt
 =X_w(b)-\max_{p\mid b}w_p,
\]
where the maximum is zero when the weighted support is empty. The
difference between this integral and the needed hinge is
\[
 \min(X_w(b),1)-\max_{p\mid b}w_p\ge0.
\]
It need not be a negligible error: when $w_2=1/2$ and all other weights
vanish, it is $1/2$ at every multiple of four, totaling about $M/8$.
Even a binary proof would therefore require additional control of these
gaps to extend by this averaging argument to arbitrary real weights.

\section{Reproduction and final status}
The accompanying archive contains this \LaTeX{} source, its PDF, Python
verification scripts, and machine-readable output. No third-party paper
sources are bundled. The standard-library checks can be run with
\begin{verbatim}
python checks/verify_exact_certificates.py
python checks/verify_hub_obstruction.py
python checks/verify_hub_family.py
python checks/verify_omega2_certificate.py
\end{verbatim}
The matching search additionally uses NumPy and SciPy. The construction
stress check is separate from the theorem's proof. The product-cover
dynamic program processes each interval position exactly once and stores
the least cardinality for each capped valuation state; its optima are
therefore exact, without floating-point optimization.

The all-scale extension and conditional transfer were independently
rederived by a separate reasoning pass, including the distinction between
the target product and incidental prime factors of selected positions.
This is a mathematical audit, not a Lean proof or human peer review.

\paragraph{What remains unresolved.}
The missing step is a proof of $(H_2)$, a sufficiently uniform weakened
version implying a $2+o(1)$ fractional coefficient, or a different direct
covering argument. Even proving $(H_2)$ and applying the rounding above
would leave the stronger exact $2n$ bound to be settled. The computations,
special classes, and counterexamples to auxiliary rules in this note do
not fill those gaps.

\begin{thebibliography}{9}
\raggedright
\bibitem{problem} T. F. Bloom (editor), Erd\H{o}s Problem 708.
\url{https://www.erdosproblems.com/708}.
\bibitem{wang} Ziming Wang, \emph{Erd\H{o}s 708 interval divisor covers},
repository and source \texttt{papers/main-paper/raw-transfer.tex}.
\url{https://github.com/science-prince/erdos-708-interval-divisor-covers}.
\bibitem{berg-asym} C. Bergstr\"om, \emph{Erd\H{o}s 708 asymptotic},
whole-cofactor peeling and hinge reductions.
\url{https://github.com/pathscale/erdos-708-asymptotic};
\url{https://doi.org/10.5281/zenodo.24342718}.
\bibitem{berg-hinge} C. Bergstr\"om,
\emph{Scaled-Interval Hinge Telescoping, Difference-of-Convex Obstructions,
and Smooth Proper-Divisor Matching for Erd\H{o}s Problem 708},
19 November 2026. \url{https://doi.org/10.5281/zenodo.24366998}.
\bibitem{chen} Haoyu Chen,
\emph{Explicit upper bounds for the Erd\H{o}s--Sur\'anyi function $g(n)$},
version 14, source repository.
\url{https://github.com/chy4pro/erdos-708-explicit-upper-bounds}.
\bibitem{preil} D. Preil, \emph{Erd\H{o}s 708 notes}, Theorem B.
\url{https://github.com/dpreil-ydlp/open-problems/blob/master/math/erdos708/NOTES.md}.
\bibitem{tripodo} John Tripodo,
\emph{Sharp finite cases of the Erd\H{o}s--Sur\'anyi interval-product
problem: $g(4)=5$ and $g(5)=6$}, manuscript dated 19 November 2026,
Hexagon version 1. \url{https://hexagonmath.org/2611.00694v1}.
\end{thebibliography}
\end{document}
